\section{Historical Situation}

Early on 2 July 1644 the Allied army began marching toward Tadcaster. Its horse
remained behind as a rearguard, but much of its infantry, artillery, and baggage
was already on the road when Rupert appeared in strength. The Allied commanders
hurriedly recalled the column and began rebuilding their line on Marston Hill.

Newcastle joined Rupert during the morning, but Eythin did not bring the York
infantry onto the moor until late afternoon. By then the opportunity had passed:
the Allied deployment was substantially complete between 2 and 3 p.m., while
Newcastle's infantry arrived between 4 and 5 p.m. The historical evidence and
timing are summarized in Historic England's
\href{https://historicengland.org.uk/content/docs/listing/battlefields/marston-moor/}{battlefield report}
and the \href{https://bcw-project.org.uk/military/english-civil-war/northern-england/battle-of-marston-moor}{BCW Project account}.

Here Eythin gets the York regiments moving promptly. Rupert has his full army
at noon and attacks while the Allied horse is still covering an infantry column
that is arriving piecemeal.

\begin{scenarioBox}{What the Scenario Is Testing}
The Royalists receive a brief superiority in formed foot and artillery, but the
advantage erodes every half-hour. The Royalist player must press the Allied
rearguard or seize the forming ground before Leven's large infantry contingent
arrives. The Allied player must trade space for time without allowing the horse
to be trapped or the assembly area to be overrun.
\end{scenarioBox}

\subsection{Length}

The scenario lasts 12 half-hour turns, beginning at noon. Advance the time by
30 minutes each turn; the scenario ends at 6 p.m. after Turn 12.
